\documentclass[10pt]{article}
\usepackage[margin=1in]{geometry}
\usepackage{amsmath,amssymb,amsthm,mathtools}
\usepackage[T1]{fontenc}
\usepackage{lmodern}
\usepackage[hidelinks]{hyperref}
\newtheorem{theorem}{Theorem}
\newtheorem{lemma}{Lemma}
\newcommand{\II}{\mathrm{II}}
\newcommand{\R}{\mathbb R}
\newcommand{\Z}{\mathbb Z}
\newcommand{\N}{\mathbb N}
\newcommand{\dimH}{\dim_{\mathrm H}}
\title{A proof of conjecture 00000000302}
\author{C0ldSmi1e\\\small Prepared with Codex assistance}
\date{}
\begin{document}
\maketitle
\begin{abstract}
For every positive exponent parameter, the exponent-badly-approximable set in the conjecture has full Lebesgue measure in the real line. Its Hausdorff dimension, and that of the intersection at parameters $1/3$ and $1/2$, is therefore one. The complete argument is formalized in Lean 4.19.0 with Mathlib v4.19.0.
\end{abstract}

\section{Statement and conventions}
The English and Chinese source define
$\II_\alpha=\{\xi:\exists c>0,\ |\xi-p/q|>c/q^{2+\alpha}\}$
and ask whether
\[
  \dimH(\II_{1/3}\cap\II_{1/2})=\dimH(\II_{1/2}).
\]
The source leaves the domains and quantifiers on $p,q$ implicit. We make the usual meaning of the named Diophantine set explicit:
\[
\II_\alpha=\left\{x\in\R:\ \exists c\in\R,\ c>0,\quad
 \forall p\in\Z,\ \forall q\in\N_{>0},\quad
 \frac{c}{q^{2+\alpha}}<\left|x-\frac pq\right|\right\}.
\]
The positive constant is uniform in both $p$ and $q$; it may depend on $x$ and $\alpha$. All positive denominators are included, with no coprimality restriction. Powers are real powers and dimension is the ordinary metric Hausdorff dimension on $\R$. The argument does not restrict $x$ to a unit interval. The exact bilingual source accompanies the submission as \texttt{ORIGINAL.md}.

\begin{theorem}\label{thm:main}
For every $\alpha>0$, $\II_\alpha$ has full Lebesgue measure and $\dimH(\II_\alpha)=1$. For every $\alpha,\beta>0$,
\[
\dimH(\II_\alpha\cap\II_\beta)=1=\dimH(\II_\beta).
\]
In particular, conjecture 00000000302 is true under the conventional quantifiers above.
\end{theorem}

\section{A full-measure approximation bound}
\begin{lemma}\label{lem:ae}
Fix $s>2$. For almost every $x\in\R$, there is no real $C$ for which infinitely many denominators $q\ge1$ admit an integer $p$ with
\[
 x\ne p/q,\qquad |x-p/q|<C/q^s.
\]
\end{lemma}
\begin{proof}
Fix $m\in\Z$ and a positive integer $K$. On $[m,m+1]$, let $E_q$ consist of the points within $K/q^s$ of some $p/q$. Since $K/q^s\le K$, only numerators in the integer interval from $(m-K)q$ to $(m+1+K)q$ contribute. There are at most $(1+2K)q+1$ of them. Thus
\[
 |E_q|\le 2K\bigl((1+2K)q+1\bigr)q^{-s},
 \qquad \sum_{q=1}^{\infty}|E_q|<\infty,
\]
because $1-s<-1$ and $-s<-1$. The first Borel--Cantelli lemma makes the infinitely-often exceptional set null. Take a countable union over $m\in\Z$ and $K\ge1$. Every real $C$ is bounded above by some such $K$, proving the claim on all of $\R$.
\end{proof}

\section{One strict constant for every denominator}
Fix $\alpha>0$ and set $s=2+\alpha$. Choose an irrational $x$ satisfying Lemma~\ref{lem:ae}; this holds almost everywhere because the rationals are countable. Applying the lemma at $C=1$ gives $N\in\N$ such that
\[
 |x-p/q|\ge q^{-s}\quad\text{for every }q\ge\max(1,N)
 \text{ and every }p\in\Z.
\]
For each fixed $q>0$, the closed lattice $q^{-1}\Z$ does not contain $x$, so its distance from $x$ is positive. There are only finitely many denominators $1\le q\le N$. Hence some $\varepsilon>0$ satisfies
\[
 |x-p/q|\ge\varepsilon\quad(1\le q\le N,\ p\in\Z).
\]
When this denominator range is empty, choose any positive $\varepsilon$. Put $c=\min(1,\varepsilon)/2>0$. For $q\ge N$, with $q>0$,
\[
 c/q^s<1/q^s\le|x-p/q|.
\]
For $1\le q<N$, we have $q^s\ge1$, and therefore
\[
 c/q^s\le c<\varepsilon\le|x-p/q|.
\]
These two cases prove the required strict inequality for all numerators and all positive denominators. Thus $\II_\alpha$ has full Lebesgue measure.

\section{Hausdorff dimension and the conjecture}
A full-measure subset $S\subseteq\R$ agrees almost everywhere with $\R$ and has infinite Lebesgue measure. In the normalization used by Mathlib, one-dimensional Hausdorff measure on $\R$ equals Lebesgue measure. It follows that $\dimH S\ge1$. Monotonicity and $\dimH\R=1$ give the reverse inequality.

The intersection of two full-measure sets has full measure: its complement is a union of two null sets. Apply the preceding paragraph first to $\II_\beta$ and then to $\II_\alpha\cap\II_\beta$. This proves Theorem~\ref{thm:main}; substituting $\alpha=1/3$ and $\beta=1/2$ proves the original equality.

\section{Formal correspondence and reproduction}
The file \texttt{lean/Conjecture302.lean} defines \texttt{Conjecture302.II} exactly as in Section~1. 

The theorem \texttt{mem\_II\_of\_not\_liouvilleWith} proves the strict uniform-constant argument. The full-measure statement is \texttt{ae\_mem\_II}; \texttt{dimH\_II} and \texttt{dimH\_inter\_II} prove the dimension results. The final theorem \texttt{Conjecture302.conjecture} has precisely the requested equality as its type, with no additional hypotheses.

The measure input is Mathlib's theorem \texttt{ae\_not\_liouvilleWith}, whose definition uses actual integer numerators, natural denominators tending to infinity, and real approximation errors. Mathlib proves it from convergent power-series bounds and Borel--Cantelli. Its bounded-denominator separation theorem for irrational numbers supplies the finite-denominator step. Dimension is Mathlib's genuine \texttt{dimH}, valued in $[0,\infty]$, and the proof obtains the exact finite value one.

The project pins Lean 4.19.0 and Mathlib commit
\begin{center}\small\texttt{c44e0c8ee63ca166450922a373c7409c5d26b00b}.\end{center}
The included audit prints all seven theorem types and their axiom dependencies. No auxiliary numerical computation is needed. Reproduction instructions, complete compiler output, independent semantic and engineering review records, and artifact hashes accompany the submission. Local verification and independent model review are separate from maintainer acceptance.
\end{document}
